% Apunts de teoria en PDF · 1r ESO · Unitat 4: Fraccions
% Mateix estil que els apunts de Batxillerat (franja i caixes de teoria en blau, exemples en roig), però més visuals,
% curts i amb dibuixos, per a l'alumnat de 3r ESO. Es compila amb pdflatex dues vegades (per als enllaços de l'índex),
% amb logo-fentmates.pdf (_eines/logo/pdf/horitzontal-clar.pdf) al costat.
\documentclass[11pt,a4paper]{article}
\usepackage[utf8]{inputenc}
\usepackage[T1]{fontenc}
\usepackage[catalan]{babel}
\usepackage[sfdefault,book,lining]{FiraSans}
\usepackage[italic,noplusnominus,noequal]{mathastext}
\linespread{1.18}
\usepackage[a4paper,top=2.4cm,bottom=2.2cm,left=1.6cm,right=1.6cm,headheight=1.1cm,headsep=0.5cm,footskip=1cm]{geometry}
\usepackage{amsmath,amssymb,cancel}
\usepackage{icomma}
\usepackage{graphicx}
\usepackage[table]{xcolor}
\usepackage{tikz}\usetikzlibrary{calc,positioning,arrows.meta,decorations.pathreplacing,shapes.geometric}
\usepackage{fontawesome5}
\usepackage{ccicons}
\usepackage{fancyhdr}
\usepackage{enumitem}
\usepackage{array}
\usepackage{needspace}
\usepackage[most]{tcolorbox}
\usepackage{titlesec}
\usepackage[hidelinks]{hyperref}
\hypersetup{pdftitle={Apunts de teoria. Unitat 4: Fraccions. 1r ESO},pdfsubject={Matemàtiques 1r ESO},pdfkeywords={fracció, numerador, denominador, fracció d'un nombre, percentatges, fraccions equivalents, amplificar, simplificar, fracció irreductible, productes creuats, comú denominador, comparar fraccions, suma, resta, multiplicació, divisió, operacions combinades, problemes}}

% colors de la web (bloc d'Àlgebra: blau)
\definecolor{fmnavy}{HTML}{1B2748}
\definecolor{fmblue}{HTML}{3D6AA0}
\definecolor{fmbluel}{HTML}{4E81BD}
\definecolor{fmyellow}{HTML}{FFC94D}
\definecolor{fmgold}{HTML}{E9A919}
\definecolor{fmaccent}{HTML}{FF7A50}
\definecolor{fmdef}{HTML}{1D4F73}\definecolor{fmdefbg}{HTML}{EEF4F9}
\definecolor{fmex}{HTML}{A3203A}\definecolor{fmexbg}{HTML}{FBF0F2}
\definecolor{fmsf}{HTML}{C2410C}
\definecolor{fmok}{HTML}{15803D}
\definecolor{fmmuted}{HTML}{5D6685}
\definecolor{fmsurface}{HTML}{EEF1F8}
\definecolor{fmA}{HTML}{2F6FD6}\definecolor{fmB}{HTML}{EA6A10}\definecolor{fmR}{HTML}{E11D48}\definecolor{fmK}{HTML}{16A34A}\definecolor{fmV}{HTML}{7C3AED}

% lletres
\newcommand{\comic}{\fontfamily{ComicNeue-TLF}\selectfont}
\newcommand{\heavy}{\fontseries{eb}\selectfont}
\newcommand{\semi}{\fontseries{sb}\selectfont}
\DeclareUnicodeCharacter{2212}{-}

\setlength{\parindent}{0pt}
\setlength{\parskip}{7pt}
\setlength{\lineskip}{4pt}\setlength{\lineskiplimit}{4pt}
\setlength{\emergencystretch}{3em}
\setlist{itemsep=3pt,topsep=3pt,leftmargin=1.4em}
\setcounter{secnumdepth}{1}
\renewcommand{\thesection}{4.\arabic{section}}
\titleformat{\section}[hang]{\color{fmnavy}\fontsize{23}{27}\heavy\selectfont}{\tcbox[on line,colback=fmblue,colframe=fmblue,arc=3mm,boxsep=2pt,left=6pt,right=6pt,top=3pt,bottom=3pt]{\color{white}\thesection}}{0.45em}{}[\vspace{5pt}{\color{fmblue}\titlerule[2pt]}]
\titlespacing*{\section}{0pt}{4pt}{6pt}
\titleformat{\subsection}{\color{fmblue}\Large\semi}{}{0pt}{\faCaretRight\ }
\titlespacing*{\subsection}{0pt}{8pt}{2pt}

% frase de cada apartat, «a mà»
\newcommand{\frase}[1]{\par\vspace{-2pt}{\comic\Large\color{fmbluel}\faCommentDots[regular]\ \ #1\par}\vspace{2pt}}
\newcommand{\seccio}[3]{\Needspace{16\baselineskip}\section{#2}\hypertarget{#1}{}\frase{#3}}
% notes a mà dins dels exemples
\newcommand{\nota}[1]{{\comic\small\color{fmex!85!black}#1}}
\newcommand{\notab}[1]{{\comic\small\color{fmblue}#1}}
% bé / malament
\newcommand{\be}[1]{{\color{fmok}\faCheckCircle}\ \ensuremath{#1}}
\newcommand{\mal}[1]{{\color{fmR}\faTimesCircle}\ {\color{fmR}\ensuremath{#1}}}
% pas numerat
\newcommand{\pas}[1]{\tikz[baseline=(p.base)]\node[circle,fill=fmblue,text=white,inner sep=1.2pt,minimum size=1.25em,font=\small\heavy](p){#1};}
% etiqueta de color
\newcommand{\chip}[2]{\tcbox[on line,colback=#1!14,colframe=#1!60,arc=2mm,boxsep=1pt,left=5pt,right=5pt,top=1pt,bottom=1pt]{#2}}
\newcommand{\fmkey}[1]{\colorbox{fmyellow!55}{\strut #1}}
\newcommand{\cA}[1]{{\color{fmA}#1}}\newcommand{\cB}[1]{{\color{fmB}#1}}\newcommand{\cK}[1]{{\color{fmK}#1}}\newcommand{\cR}[1]{{\color{fmR}#1}}\newcommand{\cV}[1]{{\color{fmV}#1}}

% caixes
\tcbset{fmbase/.style={enhanced,arc=3mm,boxrule=1.2pt,left=4mm,right=4mm,top=2mm,bottom=2.5mm,before skip=8pt,after skip=8pt,fonttitle=\heavy\large,before upper=\raggedright}}
\newtcolorbox{idea}[1][Idea clau]{fmbase,colback=fmdefbg,colframe=fmdef,coltitle=white,title={\faKey\ \ #1}}
\newtcolorbox{exemple}[1][Exemple]{fmbase,colback=fmexbg,colframe=fmex,colbacktitle=fmexbg,coltitle=fmex,title={\faPencil*\ \ #1},titlerule=0.5pt,titlerule style={fmex!40,dashed}}
\newtcolorbox{compte}[1][Compte!]{fmbase,colback=fmsf!6,colframe=fmsf,coltitle=white,fonttitle=\comic\bfseries\Large,title={\faExclamationTriangle\ \ #1}}
\newtcolorbox{truc}[1][Truc]{fmbase,colback=fmyellow!20,colframe=fmgold,colbacktitle=fmyellow,coltitle=fmnavy,fonttitle=\comic\bfseries\Large,title={\faLightbulb[regular]\ \ #1}}

% dibuix de la targeta de la unitat, en blanc
\newcommand{\fmcard}{\begin{tikzpicture}[x=0.062cm,y=-0.062cm,white,line width=3pt,line cap=round]
  \begin{scope}[transparency group,opacity=0.15]
  \node[anchor=base west,inner sep=0pt,font=\fontsize{30}{30}\selectfont] at (14,44) {$\frac34\quad\frac12$};
  \fill (190,70) -- (190,25) arc (90:-180:45) -- cycle;
  \draw (190,70) circle (45);
  \end{scope}
\end{tikzpicture}}
\newcommand{\fmbrand}{\raisebox{-0.25\height}{\includegraphics[height=0.5cm]{logo-fentmates.pdf}}}

% regla de Ruffini: #1 = a, #2 = coeficients, #3 = fila del mig (sense el primer), #4 = fila de baix, #5 = nombre de coeficients
\newcommand{\ruffini}[5]{\begin{tikzpicture}[x=1.25cm,y=0.85cm,every node/.style={font=\large}]
  \foreach \v [count=\i] in {#2} \node (t\i) at (\i,0) {$\v$};
  \foreach \v [count=\i] in {#3} {\pgfmathtruncatemacro\j{\i+1}\node (m\j) at (\j,-1) {$\v$};}
  \foreach \v [count=\i] in {#4} \node (b\i) at (\i,-2.1) {$\v$};
  \node[text=fmB] at (0,-1) {$#1$};
  \draw[line width=1.1pt,fmnavy] (0.45,0.5)--(0.45,-1.5) (0.45,-1.5)--(#5+0.5,-1.5);
  \node[draw=fmR,line width=1.3pt,rounded corners=2mm,minimum size=0.75cm] at (#5,-2.1) {};
  \draw[-{Latex[length=2.2mm]},fmK,line width=1.1pt] (t1.south) -- (b1.north);
  \pgfmathtruncatemacro\last{#5-1}
  \foreach \i in {1,...,\last} {\pgfmathtruncatemacro\j{\i+1}
    \draw[-{Latex[length=2.2mm]},fmB,line width=1.1pt,shorten >=1pt] (b\i.north east) -- (m\j.south west);
    \draw[-{Latex[length=2.2mm]},fmA!70,line width=1pt,dashed,shorten >=2pt] ([xshift=8pt]t\j.south) -- ([xshift=8pt]b\j.north);}
\end{tikzpicture}}

% capçalera i peu
\pagestyle{fancy}\fancyhf{}
\fancyhead[L]{\small\fmbrand}
\fancyhead[R]{\small\color{fmmuted}1r ESO · Unitat 4: Fraccions}
\fancyfoot[L]{\scriptsize\color{fmmuted}\href{https://creativecommons.org/licenses/by-nc-sa/4.0/deed.ca}{\ccbyncsa}\ \ FentMates · fentmates.org · CC BY-NC-SA 4.0}
\fancyfoot[R]{\small\color{fmmuted}\thepage}
\renewcommand{\headrulewidth}{0.4pt}\renewcommand{\headrule}{\color{fmmuted!40}\hrule height \headrulewidth}
\fancypagestyle{primera}{\fancyhf{}\renewcommand{\headrulewidth}{0pt}
  \fancyfoot[L]{\scriptsize\color{fmmuted}\href{https://creativecommons.org/licenses/by-nc-sa/4.0/deed.ca}{\ccbyncsa}\ \ FentMates · fentmates.org · CC BY-NC-SA 4.0}
  \fancyfoot[R]{\small\color{fmmuted}\thepage}}

% targeta de l'índex
\newcommand{\itile}[4]{\begin{tcolorbox}[enhanced,colback=white,colframe=fmblue!35,arc=3mm,boxrule=1pt,boxsep=0pt,left=6pt,right=6pt,top=5pt,bottom=5pt,before skip=0pt,after skip=7pt]
  \hyperlink{#1}{\begin{minipage}[c]{0.9cm}\centering{\color{fmblue}\LARGE #3}\end{minipage}\hspace{6pt}%
  \begin{minipage}[c]{0.78\linewidth}{\color{fmblue}\heavy #2}\\[-2pt]{\color{fmnavy}\semi\large #4}\end{minipage}}\end{tcolorbox}}

% taula amb vores arrodonides: el contingut es retalla amb la forma (sense @{} a les vores ni entre columnes)
\newsavebox\tabclipbox
\newcommand{\taulaclip}[3][3mm]{\sbox\tabclipbox{#3}\begin{tikzpicture}[baseline=(tcn.base)]
  \begin{scope}\clip[rounded corners=#1] (0,-\dp\tabclipbox) rectangle (\wd\tabclipbox,\ht\tabclipbox);
  \node[inner sep=0pt,outer sep=0pt,anchor=base west] (tcn) at (0,0) {\usebox\tabclipbox};\end{scope}
  \draw[#2,rounded corners=#1] (0,-\dp\tabclipbox) rectangle (\wd\tabclipbox,\ht\tabclipbox);
\end{tikzpicture}}
\newcommand{\taularod}[3][fmblue]{\taulaclip{#1,line width=1.2pt}{\renewcommand{\arraystretch}{1.5}\begin{tabular}{#2}#3\end{tabular}}}
% balança
\newcommand{\balanca}[3]{\begin{tikzpicture}[scale=#3,transform shape,every node/.style={anchor=center}]
  \fill[fmnavy] (-0.45,-1.9) rectangle (0.45,-1.75);
  \fill[fmnavy!80] (-0.12,-1.75) rectangle (0.12,0);
  \fill[fmgold] (0,0.05) circle (0.13);
  \draw[fmnavy,line width=2.4pt,line cap=round] (-2,0) -- (2,0);
  \foreach \s in {-1,1} {
    \draw[fmmuted,line width=0.8pt] (\s*2,0) -- (\s*2-0.75,-1.05) (\s*2,0) -- (\s*2+0.75,-1.05);
    \fill[fmyellow!70,draw=fmgold,line width=1pt] (\s*2-0.95,-1.05) arc[start angle=180,end angle=360,x radius=0.95,y radius=0.32] -- cycle;}
  \node[font=\Large,fill=fmA!15,draw=fmA,rounded corners=1.5mm,inner sep=3pt] at (-2,-0.68) {$#1$};
  \node[font=\Large,fill=fmB!15,draw=fmB,rounded corners=1.5mm,inner sep=3pt] at (2,-0.68) {$#2$};
\end{tikzpicture}}
% eixos amb quadrícula (dins d'un tikzpicture): xmin, xmax, ymin, ymax
\newcommand{\eixos}[4]{\draw[fmmuted!22,line width=0.4pt] (#1,#3) grid (#2,#4);
  \draw[-{Latex[length=2mm]},fmmuted,line width=0.8pt] (#1-0.2,0)--(#2+0.4,0) node[right,font=\small]{$x$};
  \draw[-{Latex[length=2mm]},fmmuted,line width=0.8pt] (0,#3-0.2)--(0,#4+0.4) node[above,font=\small]{$y$};}
% taula de traduccions (diccionari) amb les vores arrodonides
\newcommand{\dic}[2]{#1 & {\large$#2$}\\}
\newcommand{\dicratio}{0.62}
\newcommand{\dichead}{\begin{tabular}{@{}>{\raggedright}p{\dimexpr\dicratio\linewidth-12pt}>{\centering\arraybackslash}p{\dimexpr\linewidth-\dicratio\linewidth-12pt}@{}}\faComment[regular]\ Ho diem així\dots & \dots\ així\end{tabular}}
\newenvironment{dicbox}{\begin{tcolorbox}[enhanced,arc=3mm,boxrule=1.2pt,colframe=fmblue,colback=white,coltitle=white,fonttitle=\heavy,
  title={\dichead},left=0pt,right=0pt,top=0pt,bottom=1.5mm,boxsep=0pt,lefttitle=6pt,righttitle=6pt,before skip=0pt,after skip=0pt]\rowcolors{1}{white}{fmsurface}%
  \begin{tabular}{@{\hspace{6pt}}>{\raggedright}p{\dimexpr\dicratio\linewidth-12pt}>{\centering\arraybackslash}p{\dimexpr\linewidth-\dicratio\linewidth-12pt}@{\hspace{6pt}}}}
  {\end{tabular}\end{tcolorbox}}
% resum
\newcommand{\rtile}[3]{\begin{tcolorbox}[enhanced,colback=white,colframe=fmblue!45,arc=3mm,boxrule=1pt,left=3mm,right=3mm,top=2.5mm,bottom=2.5mm,fontupper=\large,
  title={\tcbox[on line,colback=fmblue,colframe=fmblue,arc=2mm,boxsep=0pt,left=4pt,right=4pt,top=1pt,bottom=1pt]{\color{white}\small\heavy #1}\ \ #2},coltitle=fmnavy,colbacktitle=fmdefbg,fonttitle=\semi,toptitle=0.5mm,bottomtitle=0.5mm]#3\end{tcolorbox}}
% portada: #1 subtítol a mà
\newcommand{\portada}[1]{\thispagestyle{primera}
\begin{tikzpicture}[remember picture,overlay]
  \fill[fmblue] (current page.north west) rectangle ([yshift=-9.4cm]current page.north east);
  \begin{scope}
    \clip (current page.north west) rectangle ([yshift=-9.4cm]current page.north east);
    \node[anchor=north west,inner sep=0pt] at ([xshift=6.6cm,yshift=-1.3cm]current page.north west) {\fmcard};
  \end{scope}
  \node[anchor=north west,inner sep=0pt] at ([xshift=1.6cm,yshift=-1.1cm]current page.north west) {\tcbox[colback=white,colframe=white,arc=2.5mm,boxsep=0pt,left=7pt,right=7pt,top=5pt,bottom=5pt]{\includegraphics[height=0.75cm]{logo-fentmates.pdf}}};
\end{tikzpicture}
\vspace*{0.55cm}
\tcbox[colback=white!15!fmblue,colframe=white!15!fmblue,arc=2.5mm,boxsep=0pt,left=6pt,right=6pt,top=3pt,bottom=3pt]{\color{white}\small\bfseries Apunts de teoria · 1r ESO · Matemàtiques}

\vspace{2pt}
{\raggedright\color{white}\fontsize{30}{36}\heavy\selectfont{\color{fmyellow}Unitat 4}\\[2pt]Fraccions\par}
\vspace{8pt}
{\color{white}\comic\Large #1\par}
\vspace*{\dimexpr 7.4cm-\pagetotal\relax}
}
\newcommand{\llegenda}{\par\vspace{16pt}
{\color{fmnavy}\heavy\Large\faPalette\ \ Com es llegeixen aquests apunts}\par\smallskip
\begin{tcbraster}[raster columns=4,raster equal height,raster column skip=8pt,raster before skip=12pt]
\tcbox[colback=fmdefbg,colframe=fmdef,arc=3mm,boxrule=1.2pt,halign=center,width=\linewidth]{{\color{fmdef}\Large\faKey}\\[2pt]{\color{fmdef}\heavy Idea clau}\\[-1pt]{\small el que cal saber}}
\tcbox[colback=fmexbg,colframe=fmex,arc=3mm,boxrule=1.2pt,halign=center,width=\linewidth]{{\color{fmex}\Large\faPencil*}\\[2pt]{\color{fmex}\heavy Exemple}\\[-1pt]{\small pas a pas}}
\tcbox[colback=fmsf!6,colframe=fmsf,arc=3mm,boxrule=1.2pt,halign=center,width=\linewidth]{{\color{fmsf}\Large\faExclamationTriangle}\\[2pt]{\color{fmsf}\heavy Compte!}\\[-1pt]{\small errors típics}}
\tcbox[colback=fmyellow!20,colframe=fmgold,arc=3mm,boxrule=1.2pt,halign=center,width=\linewidth]{{\color{fmgold!80!black}\Large\faLightbulb[regular]}\\[2pt]{\color{fmnavy}\heavy Truc}\\[-1pt]{\small per a recordar-ho}}
\end{tcbraster}}
\newcommand{\resumtitol}{\clearpage
\hypertarget{resum}{}
{\color{fmnavy}\fontsize{23}{27}\heavy\selectfont\faClipboardList\ \ En resum}\par\vspace{5pt}{\color{fmblue}\titlerule[2pt]}
\frase{Tota la unitat en una pàgina.}}
% els 5 passos per a resoldre problemes: #2 = text del pas 2
\newcommand{\passosproblema}[2]{\begin{center}
\begin{tikzpicture}[s/.style={draw=fmblue,fill=fmdefbg,rounded corners=3mm,line width=1.2pt,minimum width=2.75cm,minimum height=2.45cm,align=center,font=\small},
  ic/.style={font=\LARGE,text=fmblue}]
  \foreach \i/\ic/\t in {1/\faBookOpen/{\textbf{Llegim}\\i entenem},2/\faQuestionCircle/{#1},3/\faEquals/{#2},4/\faCogs/{\textbf{Resolem}},5/\faCheckDouble/{\textbf{Comprovem}\\i responem}} {
    \node[s] (s\i) at (3.3*\i,0) {};
    \node[ic] at ([yshift=-0.5cm]s\i.north) {\ic};
    \node[align=center,font=\small,anchor=north] at ([yshift=-0.95cm]s\i.north) {\t};
    \node[circle,fill=fmblue,text=white,font=\small\heavy,inner sep=1.5pt,minimum size=1.3em] at (s\i.north west) {\i};}
  \foreach \i [evaluate=\i as \j using int(\i+1)] in {1,...,4} \draw[-{Latex[length=3mm]},fmmuted,line width=1.5pt] (s\i) -- (s\j);
\end{tikzpicture}
\end{center}}

\newcommand{\lo}{\mathopen{]}}\newcommand{\ro}{\mathclose{[}}
% punts ple i buit, nombres dels eixos i taula de valors
\newcommand{\pple}[2][fmblue]{\fill[#1] #2 circle (3.6pt);}
\newcommand{\nums}[4]{\foreach \k in {#1,...,#2} {\ifnum\k=0\else\node[below,font=\scriptsize,text=fmmuted] at (\k,0) {$\k$};\fi}
  \foreach \k in {#3,...,#4} {\ifnum\k=0\else\node[left,font=\scriptsize,text=fmmuted] at (0,\k) {$\k$};\fi}}
\newcommand{\tv}[3][fmblue]{\taulaclip{#1,line width=1.2pt}{\renewcommand{\arraystretch}{1.45}%
  \begin{tabular}{>{\columncolor{#1}\color{white}\heavy}c*{12}{w{c}{0.8cm}}}#2\\\rowcolor{#1!8}\cellcolor{#1}#3\end{tabular}}}
% etiqueta d'una recta o corba
\newcommand{\etq}[3]{\node[fill=white,inner sep=1.5pt,text=#1,font=\small] at #2 {#3};}

% interval sobre la recta: #1 esquerre, #2 dret, #3 tipus esquerre (o/c/a), #4 tipus dret (o/c/a), #5 color
\def\intvdot#1#2#3{\ifx#2o\filldraw[fill=white,draw=#3,line width=1.3pt] (#1,0) circle (3.3pt);\fi\ifx#2c\fill[#3] (#1,0) circle (3.5pt);\fi}
\newcommand{\intv}[5]{\begin{tikzpicture}[x=0.4cm,baseline=-3pt]
  \draw[fmmuted,line width=0.8pt] (-6,0)--(7,0);
  \foreach \k in {-5,...,6} \draw[fmmuted,line width=0.6pt] (\k,-0.08)--(\k,0.08);
  \ifx#3a\def\xa{-6}\else\def\xa{#1}\fi
  \ifx#4a\def\xb{7}\else\def\xb{#2}\fi
  \ifx#4a\draw[#5,line width=2.6pt,-{Latex[length=2.6mm]}] (\xa,0)--(\xb,0);\else\ifx#3a\draw[#5,line width=2.6pt,{Latex[length=2.6mm]}-] (\xa,0)--(\xb,0);\else\draw[#5,line width=2.6pt] (\xa,0)--(\xb,0);\fi\fi
  \intvdot{#1}{#3}{#5}\intvdot{#2}{#4}{#5}
  \ifx#3a\else\node[below=2pt,font=\scriptsize] at (#1,0) {$#1$};\fi
  \ifx#4a\else\node[below=2pt,font=\scriptsize] at (#2,0) {$#2$};\fi
\end{tikzpicture}}

\newcommand{\passos}[1]{\begin{tcolorbox}[fmbase,colback=fmdefbg,colframe=fmdef,coltitle=white,title={\faListOl\ \ Com es resolen}]#1\end{tcolorbox}}
\newcommand{\recta}[2]{\draw[-{Latex[length=2.5mm]},fmnavy,line width=1.1pt] (#1-0.6,0) -- (#2+0.6,0);
  \foreach \k in {#1,...,#2} {\draw[fmnavy,line width=0.9pt] (\k,-0.12) -- (\k,0.12); \node[below=3pt,font=\small] at (\k,0) {$\k$};}}
\newcommand{\barra}[4]{\begin{tikzpicture}
  \foreach \i in {1,...,#1} {\pgfmathsetmacro\xa{(\i-1)*#4/#1}\pgfmathsetmacro\xb{\i*#4/#1}
    \ifnum\i>#2 \filldraw[fill=white,draw=#3,line width=1pt] (\xa,0) rectangle (\xb,0.55);
    \else \filldraw[fill=#3!35,draw=#3,line width=1pt] (\xa,0) rectangle (\xb,0.55); \fi}
\end{tikzpicture}}
% pastís: #1 parts, #2 pintades, #3 color, #4 radi
\newcommand{\pastis}[4]{\begin{tikzpicture}\foreach \i in {1,...,#2} \fill[#3!45] (0,0) -- ({90-360*(\i-1)/#1}:#4) arc ({90-360*(\i-1)/#1}:{90-360*\i/#1}:#4) -- cycle;
  \foreach \i in {1,...,#1} \draw[#3,line width=0.9pt] (0,0) -- ({90-360*\i/#1}:#4); \draw[#3,line width=1.2pt] (0,0) circle (#4);\end{tikzpicture}}
% barra: #1 parts, #2 pintades, #3 color, #4 llarg
\newcommand{\barrafr}[4]{\begin{tikzpicture}\foreach \i in {1,...,#2} \fill[#3!45] ({(\i-1)*#4/#1},0) rectangle ({\i*#4/#1},0.45);
  \foreach \i in {0,...,#1} \draw[#3,line width=0.9pt] ({\i*#4/#1},0) -- ({\i*#4/#1},0.45); \draw[#3,line width=1.2pt] (0,0) rectangle (#4,0.45);\end{tikzpicture}}


\begin{document}
\portada{Partim la unitat en parts iguals i en prenem unes quantes.}

{\color{fmnavy}\heavy\Large\faMapSigns\ \ Què aprendràs?}\par\smallskip
\begin{minipage}[t]{0.485\linewidth}
\itile{sec1}{4.1}{\faChartPie}{El concepte de fracció}
\itile{sec2}{4.2}{\faIcon{percentage}}{Fracció d'un nombre}
\itile{sec3}{4.3}{\faPercent}{Percentatges i fraccions}
\itile{sec4}{4.4}{\faEquals}{Fraccions equivalents}
\itile{sec5}{4.5}{\faTimes}{Són equivalents?}
\itile{sec6}{4.6}{\faLayerGroup}{Comú denominador}
\end{minipage}\hfill
\begin{minipage}[t]{0.485\linewidth}
\itile{sec7}{4.7}{\faSortAmountUp}{Comparar i ordenar}
\itile{sec8}{4.8}{\faPlus}{Sumes i restes}
\itile{sec9}{4.9}{\faDivide}{Multiplicació i divisió}
\itile{sec10}{4.10}{\faCalculator}{Operacions combinades}
\itile{sec11}{4.11}{\faLightbulb}{Problemes amb fraccions}
\itile{resum}{\phantom{4.11}}{\faClipboardList}{En resum}
\end{minipage}
\llegenda

% ============================================================ 4.1
\clearpage
\seccio{sec1}{El concepte de fracció}{Partim la unitat en parts iguals i en prenem unes quantes.}

\vfill
\begin{minipage}[c]{0.55\linewidth}
\begin{idea}[Què és una fracció?]
{\centering\Large$\dfrac{a}{b}$ \ \ \normalsize amb $b\neq0$\par}\smallskip
El \cA{\textbf{denominador}} $b$ (baix): en quantes parts \textbf{iguals} dividim la unitat.\\[2pt]
El \cB{\textbf{numerador}} $a$ (dalt): quantes d'aquestes parts agafem.
\end{idea}
\end{minipage}\hfill
\begin{minipage}[c]{0.42\linewidth}\centering
\pastis{4}{3}{fmA}{0.75} \qquad \pastis{8}{3}{fmB}{0.75}\\[4pt]
{\large$\frac34$ \qquad\qquad\ $\frac38$}\\[6pt]
\barrafr{5}{2}{fmK}{4}\\[2pt]{\large$\frac25$}
\end{minipage}

\vfill
\begin{compte}[Les parts han de ser iguals]
Si partim una pizza en 4 trossos de mides diferents, cada tros \textbf{no} és $\frac14$.
\end{compte}

\vfill
\begin{tcbraster}[raster columns=3,raster equal height,raster column skip=8pt]
\begin{tcolorbox}[fmbase,colback=fmA!8,colframe=fmA,coltitle=white,title={Part de la unitat}]
$\frac34$ d'una pizza: la dividim en 4 parts iguals i en mengem 3.
\end{tcolorbox}
\begin{tcolorbox}[fmbase,colback=fmB!8,colframe=fmB,coltitle=white,title={Una divisió}]
$\frac34=3:4=0{,}75$. 3 pizzes entre 4 amics. \ $\frac{12}4=3$: un natural.
\end{tcolorbox}
\begin{tcolorbox}[fmbase,colback=fmK!8,colframe=fmK,coltitle=white,title={Un operador}]
Actua sobre una quantitat: $\frac34$ de 20 és 15 (4.2).
\end{tcolorbox}
\end{tcbraster}

\vfill
\begin{minipage}[c]{0.5\linewidth}\centering\small
\taularod{>{\hspace{3pt}}c<{\hspace{3pt}}>{\hspace{3pt}}l<{\hspace{3pt}}>{\hspace{3pt}}c<{\hspace{3pt}}>{\hspace{3pt}}l<{\hspace{3pt}}}{
\rowcolor{fmblue}\multicolumn{4}{c}{\color{white}\heavy Com es llegeixen}\\
$\frac12$ & un mig & $\frac16$ & un sisé\\
\rowcolor{fmsurface}$\frac23$ & dos terços & $\frac37$ & tres setens\\
$\frac34$ & tres quarts & $\frac38$ & tres huitens\\
\rowcolor{fmsurface}$\frac45$ & quatre cinquens & $\frac{5}{12}$ & cinc dotzens
}
\end{minipage}\hfill
\begin{minipage}[c]{0.47\linewidth}
\begin{idea}[Pròpies i impròpies]
\textbf{Pròpia}: numerador menor, val \textbf{menys d'1}: $\frac34<1$.\\[2pt]
Iguals: val \textbf{1}: $\frac55=1$.\\[2pt]
\textbf{Impròpia}: numerador major, val \textbf{més d'1}: $\frac74=1+\frac34$ \nota{$7:4$, quocient 1 i residu 3}
\end{idea}
\end{minipage}

\vfill
% ============================================================ 4.2
\clearpage
\seccio{sec2}{Fracció d'un nombre}{Dividim entre el denominador i multipliquem pel numerador.}

\vfill
\begin{tcolorbox}[enhanced,colback=fmyellow!18,colframe=fmgold,arc=4mm,boxrule=1.6pt,before upper=\centering,top=2mm,bottom=2mm]
{\Large$\dfrac35\ \text{de}\ 40=(40:5)\cdot3=8\cdot3=\boxed{24}$}\par\smallskip
{\small o bé, perquè «de» vol dir multiplicar: \ $\dfrac{3\cdot40}{5}=\dfrac{120}{5}=24$}
\end{tcolorbox}

\vfill
\begin{center}
\begin{tikzpicture}[x=2cm]
  \foreach \i in {0,...,4} {\filldraw[fill=\ifnum\i<3 fmB!40\else white\fi,draw=fmB,line width=1.1pt] (\i,0) rectangle (\i+1,0.7); \node at (\i+0.5,0.35) {8};}
  \draw[decorate,decoration={brace,amplitude=5pt},line width=0.9pt] (0,0.85) -- (5,0.85) node[midway,above=6pt] {40};
  \draw[decorate,decoration={brace,amplitude=5pt,mirror},line width=0.9pt,fmB] (0,-0.15) -- (3,-0.15) node[midway,below=6pt,text=fmB] {$\frac35$ de 40 $=24$};
\end{tikzpicture}
\end{center}

\vfill
\begin{idea}[El problema invers: coneixem la part, calculem el total]
Si $\frac35$ d'un nombre són 24, cada cinqué val $24:3=8$ i el total són 5 cinquens: $8\cdot5=\boxed{40}$.\\[3pt]
{\centering$\text{Total}=(\text{part}:\text{numerador})\cdot\text{denominador}$\par}
\end{idea}

\vfill
\begin{exemple}[Exemples]
\begin{minipage}[t]{0.48\linewidth}
$\frac23$ de 60: \ $(60:3)\cdot2=\boxed{40}$\\[4pt]
$\frac34$ de 200 €: \ $(200:4)\cdot3=\boxed{150\text{ €}}$\\[4pt]
$\frac{5}{6}$ de 1 hora: \ $(60:6)\cdot5=\boxed{50\text{ min}}$
\end{minipage}\hfill
\begin{minipage}[t]{0.49\linewidth}
$\frac27$ d'un nombre són 18: \ $(18:2)\cdot7=\boxed{63}$\\[4pt]
$\frac34$ dels alumnes són 21: \ $(21:3)\cdot4=\boxed{28}$ alumnes\\[4pt]
$\frac{1}{5}$ d'un dipòsit són 12 L: \ $12\cdot5=\boxed{60}$ L
\end{minipage}
\end{exemple}

\vfill
\begin{exemple}[Exemple: \ en una classe de 28 alumnes, $\frac37$ fan esport. Quants en fan? I quants no?]
$\frac37$ de 28: \ $(28:7)\cdot3=4\cdot3=\boxed{12}$ fan esport \qquad No en fan: \ $28-12=\boxed{16}$ \ \nota{o $\frac47$ de 28}
\end{exemple}

\vfill
\begin{truc}[Fes un dibuix]
Una barra partida en tantes parts com diu el denominador ajuda molt: primer busquem quant val \textbf{una} part.
\end{truc}

\vfill
% ============================================================ 4.3
\clearpage
\seccio{sec3}{Percentatges i fraccions}{Un percentatge és una fracció de denominador 100.}

\vfill
\begin{minipage}[c]{0.45\linewidth}
\begin{idea}[Percentatge]
El 35\,\% («trenta-cinc per cent») vol dir 35 de cada 100:\\[4pt]
{\centering\Large$35\,\%=\dfrac{35}{100}$\par}
\end{idea}
\end{minipage}\hfill
\begin{minipage}[c]{0.52\linewidth}\centering\small
\taularod{>{\heavy}c<{\hspace{4pt}}>{\hspace{4pt}}c<{\hspace{4pt}}>{\hspace{4pt}}l<{\hspace{4pt}}}{
\rowcolor{fmblue}\color{white}\% & \color{white}\heavy Fracció & \color{white}\heavy És\dots\\
50\,\% & $\frac{50}{100}=\frac12$ & la meitat\\
\rowcolor{fmsurface}25\,\% & $\frac{25}{100}=\frac14$ & la quarta part\\
75\,\% & $\frac{75}{100}=\frac34$ & tres quarts\\
\rowcolor{fmsurface}20\,\% & $\frac{20}{100}=\frac15$ & la cinquena part\\
10\,\% & $\frac{10}{100}=\frac1{10}$ & la desena part\\
\rowcolor{fmsurface}100\,\% & $1$ & tot
}
\end{minipage}

\vfill
\begin{tcbraster}[raster columns=2,raster equal height,raster column skip=10pt]
\begin{tcolorbox}[fmbase,colback=fmK!8,colframe=fmK,coltitle=white,title={Percentatge d'un nombre}]
Multipliquem pel percentatge i dividim entre 100:\par\smallskip
{\centering$20\,\%$ de $180=\dfrac{20\cdot180}{100}=\boxed{36}$\par}
\end{tcolorbox}
\begin{tcolorbox}[fmbase,colback=fmA!8,colframe=fmA,coltitle=white,title={De fracció a percentatge}]
Una fracció equivalent amb denominador 100:\par\smallskip
{\centering$\dfrac35=\dfrac{3\cdot20}{5\cdot20}=\dfrac{60}{100}=\boxed{60\,\%}$\par}
\end{tcolorbox}
\end{tcbraster}

\vfill
\begin{minipage}[c]{0.36\linewidth}\centering
\begin{tikzpicture}[scale=0.32]
  \foreach \i in {0,...,9} \foreach \j in {0,...,9} {\pgfmathtruncatemacro\n{\j*10+\i}\ifnum\n<35 \fill[fmB!45] (\i,\j) rectangle ++(1,1);\fi}
  \draw[fmmuted,step=1] (0,0) grid (10,10); \draw[fmnavy,line width=1pt] (0,0) rectangle (10,10);
\end{tikzpicture}\\
{\comic\small\color{fmmuted}35 quadrets de 100: el 35\,\%}
\end{minipage}\hfill
\begin{minipage}[c]{0.6\linewidth}
\begin{truc}[Càlcul mental]
10\,\%: dividir entre 10. \ 50\,\%: la meitat. \ 25\,\%: la quarta part.\\[2pt]
El 30\,\% de 70: el 10\,\% és 7, el 30\,\% és $3\cdot7=21$.
\end{truc}
\begin{exemple}[Exemples]
25\,\% de 48: \ $48:4=\boxed{12}$ \qquad 10\,\% de 350: \ $\boxed{35}$\\[3pt]
15\,\% de 60: \ $\frac{15\cdot60}{100}=\boxed{9}$ \qquad $\frac{7}{20}=\frac{35}{100}=\boxed{35\,\%}$
\end{exemple}
\end{minipage}

\vfill
\begin{exemple}[Exemple: \ unes sabatilles de 80 € tenen un 25\,\% de descompte. Quant paguem?]
Descompte: \ 25\,\% de 80 $=80:4=20$ € \qquad Paguem: \ $80-20=\boxed{60\text{ €}}$ \ \nota{o el 75\,\% de 80}
\end{exemple}

\vfill
% ============================================================ 4.4
\clearpage
\seccio{sec4}{Fraccions equivalents}{La mateixa quantitat, escrita amb nombres diferents.}

\vfill
\begin{minipage}[c]{0.5\linewidth}\centering
\barrafr{2}{1}{fmA}{5}\\[2pt]\barrafr{4}{2}{fmB}{5}\\[2pt]\barrafr{8}{4}{fmK}{5}\\[4pt]
{\Large$\dfrac12=\dfrac24=\dfrac48$}
\end{minipage}\hfill
\begin{minipage}[c]{0.47\linewidth}
\begin{idea}[Fraccions equivalents]
Representen la \textbf{mateixa quantitat}, encara que s'escriguen amb nombres diferents.
\end{idea}
\end{minipage}

\vfill
\begin{tcbraster}[raster columns=2,raster equal height,raster column skip=10pt]
\begin{tcolorbox}[fmbase,colback=fmdefbg,colframe=fmdef,coltitle=white,title={\faSearchPlus\ \ Amplificar}]
Multipliquem numerador i denominador pel \textbf{mateix} nombre:\par\smallskip
{\centering$\dfrac23=\dfrac{2\cdot4}{3\cdot4}=\dfrac{8}{12}$\par}
\end{tcolorbox}
\begin{tcolorbox}[fmbase,colback=fmdefbg,colframe=fmdef,coltitle=white,title={\faSearchMinus\ \ Simplificar}]
Dividim tots dos entre un \textbf{divisor comú}:\par\smallskip
{\centering$\dfrac{12}{18}=\dfrac{6}{9}=\dfrac{2}{3}$ \ \nota{entre 2 i entre 3}\par}
\end{tcolorbox}
\end{tcbraster}

\vfill
\begin{idea}[Fracció irreductible]
Ja no es pot simplificar: numerador i denominador només tenen l'1 com a divisor comú. D'un sol pas: dividim entre el \textbf{MCD} (unitat 2): \ $\dfrac{12}{18}=\dfrac{12:6}{18:6}=\dfrac23$.
\end{idea}

\vfill
\begin{exemple}[Exemples]
\begin{minipage}[t]{0.48\linewidth}
Tres equivalents a $\frac35$: \ $\frac6{10}$, $\frac9{15}$, $\frac{12}{20}$\\[4pt]
$\frac{4}{7}$ amb denominador 35: \ $\frac{4\cdot5}{7\cdot5}=\boxed{\frac{20}{35}}$
\end{minipage}\hfill
\begin{minipage}[t]{0.49\linewidth}
Irreductible: \ $\frac{30}{45}=\frac{30:15}{45:15}=\boxed{\frac23}$\\[4pt]
$\frac{24}{36}=\frac{12}{18}=\frac{2}{3}$ \qquad $\frac{40}{100}=\frac25$
\end{minipage}
\end{exemple}

\vfill
\begin{truc}[Infinites, però només una d'irreductible]
Una fracció té infinites fraccions equivalents, però només una d'irreductible. Dona sempre els resultats \textbf{simplificats}.
\end{truc}

\vfill
% ============================================================ 4.5
\clearpage
\seccio{sec5}{Comprovar si són equivalents}{Els productes creuats han de ser iguals.}

\vfill
\begin{tcolorbox}[enhanced,colback=fmyellow!18,colframe=fmgold,arc=4mm,boxrule=1.6pt,before upper=\centering,top=2mm,bottom=2mm]
{\Large$\dfrac ab=\dfrac cd\iff a\cdot d=b\cdot c$}\par\smallskip
{\small multipliquem en creu}
\end{tcolorbox}

\vfill
\begin{minipage}[c]{0.4\linewidth}\centering
\begin{tikzpicture}[font=\Large]
  \node (a) at (0,0.5) {$8$}; \node (b) at (0,-0.5) {$32$}; \draw (-0.35,0) -- (0.35,0);
  \node (c) at (1.6,0.5) {$3$}; \node (d) at (1.6,-0.5) {$12$}; \draw (1.25,0) -- (1.95,0);
  \draw[fmR,line width=1.2pt,-{Latex}] (a) -- (d); \draw[fmA,line width=1.2pt,-{Latex}] (b) -- (c);
\end{tikzpicture}
\end{minipage}\hfill
\begin{minipage}[c]{0.57\linewidth}
\begin{exemple}[Exemple: \ són equivalents $\frac{8}{32}$ i $\frac{3}{12}$?]
$\cR{8\cdot12=96}$ \qquad $\cA{32\cdot3=96}$ \qquad \be{} \ sí\\[3pt]
$\frac{2}{3}$ i $\frac{5}{8}$: \ $16\neq15$ \ \mal{\text{no}}
\end{exemple}
\end{minipage}

\vfill
\begin{idea}[Calcular un terme desconegut]
Igualem els productes creuats:\\[3pt]
{\centering$\dfrac{3}{14}=\dfrac{x}{42}\ \Rightarrow\ 14\cdot x=3\cdot42=126\ \Rightarrow\ x=126:14=\boxed{9}$\par}
\end{idea}

\vfill
\begin{truc}[De cap]
$42=14\cdot3$: el numerador també es multiplica per 3, \ $x=3\cdot3=9$.
\end{truc}

\vfill
\begin{exemple}[Exemples]
\begin{minipage}[t]{0.48\linewidth}
$\frac{6}{9}$ i $\frac{10}{15}$: \ $90=90$ \ \be{}\\[4pt]
$\frac{4}{6}$ i $\frac{6}{10}$: \ $40\neq36$ \ \mal{\text{no}}\\[4pt]
$\frac{5}{x}=\frac{15}{21}$: \ $15x=105$, \ \boxed{x=7}
\end{minipage}\hfill
\begin{minipage}[t]{0.49\linewidth}
$\frac{x}{8}=\frac{9}{24}$: \ $24x=72$, \ \boxed{x=3}\\[4pt]
$\frac{2}{5}=\frac{12}{x}$: \ $2x=60$, \ \boxed{x=30}\\[4pt]
$\frac{7}{4}=\frac{x}{20}$: \ $\frac{7\cdot5}{4\cdot5}$, \ \boxed{x=35}
\end{minipage}
\end{exemple}

\vfill
\begin{compte}[En creu, no en paral·lel]
\be{8\cdot12} i \be{32\cdot3}: el numerador d'una per el denominador de l'altra. \ \mal{8\cdot3} i \mal{32\cdot12} no serveixen.
\end{compte}

\vfill
% ============================================================ 4.6
\clearpage
\seccio{sec6}{Reducció a comú denominador}{Fraccions equivalents amb el mateix denominador.}

\vfill
\begin{minipage}[c]{0.52\linewidth}
\passos{\pas{1} Calculem el \textbf{mcm} dels denominadors: serà el nou denominador.\\[3pt]
\pas{2} Per a cada fracció, dividim el mcm entre el seu denominador\dots\\[3pt]
\pas{3} \dots\ i multipliquem el numerador pel resultat.}
\end{minipage}\hfill
\begin{minipage}[c]{0.45\linewidth}
\begin{truc}[Per què el mcm?]
També serviria el producte dels denominadors, però amb el mcm els nombres són més \textbf{xicotets}.
\end{truc}
\end{minipage}

\vfill
\begin{exemple}[Exemple: \ $\frac56,\ \frac34,\ \frac7{10}$ \qquad mcm$(6,4,10)=60$]
\centering\small
\taularod[fmex]{>{\heavy}l*{3}{>{\hspace{8pt}}c<{\hspace{8pt}}}}{
\rowcolor{fmex}\color{white}Fracció & \color{white}$\frac56$ & \color{white}$\frac34$ & \color{white}$\frac7{10}$\\
mcm : denominador & $60:6=10$ & $60:4=15$ & $60:10=6$\\
\rowcolor{fmexbg}Equivalent & $\frac{5\cdot10}{60}=\frac{50}{60}$\rule[-8pt]{0pt}{22pt} & $\frac{3\cdot15}{60}=\frac{45}{60}$ & $\frac{7\cdot6}{60}=\frac{42}{60}$
}
\end{exemple}

\vfill
\begin{exemple}[Exemples]
\begin{minipage}[t]{0.48\linewidth}
$\frac12,\ \frac23$: \ mcm 6 \ $\to$ \ $\frac36,\ \frac46$\\[4pt]
$\frac34,\ \frac58$: \ mcm 8 \ $\to$ \ $\frac68,\ \frac58$
\end{minipage}\hfill
\begin{minipage}[t]{0.49\linewidth}
$\frac16,\ \frac49,\ \frac5{12}$: \ mcm 36 \ $\to$ \ $\frac6{36},\ \frac{16}{36},\ \frac{15}{36}$\\[4pt]
$3,\ \frac45$: \ $3=\frac31=\frac{15}5$ \ $\to$ \ $\frac{15}5,\ \frac45$
\end{minipage}
\end{exemple}

\vfill
\begin{compte}[Un nombre natural té denominador 1]
$5=\dfrac51$: per a passar-lo a denominador 6, \ $\dfrac{5\cdot6}{6}=\dfrac{30}{6}$.
\end{compte}

\vfill
\begin{idea}[Per a què serveix?]
Per a \textbf{comparar} fraccions (4.7) i per a \textbf{sumar-les i restar-les} (4.8): amb el mateix denominador, només hem de mirar els numeradors.
\end{idea}

\vfill
% ============================================================ 4.7
\clearpage
\seccio{sec7}{Comparar i ordenar fraccions}{Amb el mateix denominador, mirem els numeradors.}

\vfill
\begin{tcbraster}[raster columns=3,raster equal height,raster column skip=8pt]
\begin{tcolorbox}[fmbase,colback=fmA!8,colframe=fmA,coltitle=white,title={Mateix denominador},before upper=\centering]
Més gran la de \textbf{numerador} més gran.\par\smallskip
$\dfrac37<\dfrac57$
\end{tcolorbox}
\begin{tcolorbox}[fmbase,colback=fmB!8,colframe=fmB,coltitle=white,title={Mateix numerador},before upper=\centering]
Més gran la de \textbf{denominador} més xicotet: parts més grans.\par\smallskip
$\dfrac34>\dfrac37$
\end{tcolorbox}
\begin{tcolorbox}[fmbase,colback=fmK!8,colframe=fmK,coltitle=white,title={Tots dos diferents},before upper=\centering]
Comú denominador i comparem numeradors.\par\smallskip
$\dfrac25\ ?\ \dfrac56$: \ $\dfrac{12}{30}<\dfrac{25}{30}$
\end{tcolorbox}
\end{tcbraster}

\vfill
\begin{minipage}[c]{0.45\linewidth}\centering
\pastis{4}{3}{fmB}{0.8} \qquad \pastis{7}{3}{fmB}{0.8}\\[3pt]
{\large$\frac34$ \quad $>$ \quad $\frac37$}\\[2pt]
{\comic\small\color{fmmuted}els quarts són més grans que els setens}
\end{minipage}\hfill
\begin{minipage}[c]{0.52\linewidth}
\begin{truc}[Compara amb la unitat]
$\frac89$ és menor que 1 i $\frac{11}9$ és major: \ $\frac89<\frac{11}9$.\\[3pt]
Una fracció \textbf{impròpia} sempre és major que una \textbf{pròpia}.
\end{truc}
\end{minipage}

\vfill
\begin{exemple}[Exemple: \ ordena de menor a major \ $\frac35,\ \frac27,\ \frac{11}2$]
mcm$(5,7,2)=70$: \quad $\frac{42}{70},\ \frac{20}{70},\ \frac{385}{70}$ \qquad Pels numeradors: \quad $\boxed{\frac27<\frac35<\frac{11}2}$
\end{exemple}

\vfill
\begin{exemple}[Exemples]
\begin{minipage}[t]{0.48\linewidth}
$\frac{5}{8}\ ?\ \frac{3}{4}$: \ $\frac58<\frac68$ \ \ $\to$ \ $\frac58<\frac34$\\[4pt]
$\frac{2}{3}\ ?\ \frac{3}{5}$: \ $\frac{10}{15}>\frac{9}{15}$ \ $\to$ \ $\frac23>\frac35$
\end{minipage}\hfill
\begin{minipage}[t]{0.49\linewidth}
De major a menor: \ $\frac56,\ \frac34,\ \frac23$\\ \hspace*{1em}($\frac{10}{12}>\frac9{12}>\frac8{12}$)\\[4pt]
Qui ha menjat més pizza, $\frac38$ o $\frac25$? \ $\frac{15}{40}<\frac{16}{40}$: \ $\frac25$
\end{minipage}
\end{exemple}

\vfill
% ============================================================ 4.8
\clearpage
\seccio{sec8}{Sumes i restes de fraccions}{Cal que tinguen el mateix denominador.}

\vfill
\begin{tcbraster}[raster columns=2,raster equal height,raster column skip=10pt]
\begin{tcolorbox}[fmbase,colback=fmK!8,colframe=fmK,coltitle=white,title={Mateix denominador},before upper=\centering]
Sumem (o restem) els \textbf{numeradors} i deixem el denominador.\par\smallskip
$\dfrac47+\dfrac87=\dfrac{12}7$ \qquad $\dfrac59-\dfrac29=\dfrac39=\dfrac13$
\end{tcolorbox}
\begin{tcolorbox}[fmbase,colback=fmA!8,colframe=fmA,coltitle=white,title={Denominadors diferents},before upper=\centering]
Primer, \textbf{comú denominador}; després, com abans.\par\smallskip
$\dfrac23+\dfrac15=\dfrac{10}{15}+\dfrac3{15}=\dfrac{13}{15}$
\end{tcolorbox}
\end{tcbraster}

\vfill
\begin{minipage}[c]{0.5\linewidth}\centering
\barrafr{3}{2}{fmA}{4.5}\\[-2pt]{\small$\frac23=\frac{10}{15}$}\\[3pt]
\barrafr{5}{1}{fmB}{4.5}\\[-2pt]{\small$\frac15=\frac3{15}$}\\[3pt]
\barrafr{15}{13}{fmK}{4.5}\\[-2pt]{\small$\frac{13}{15}$}
\end{minipage}\hfill
\begin{minipage}[c]{0.47\linewidth}
\begin{compte}[Error típic]
Mai no sumem numeradors i denominadors per separat:\\[3pt]
\mal{\dfrac12+\dfrac13=\dfrac25} \qquad \be{\dfrac36+\dfrac26=\dfrac56}
\end{compte}
\end{minipage}

\vfill
\begin{exemple}[Exemples]
\begin{minipage}[t]{0.48\linewidth}
Amb un natural: \ $4+\frac{25}{9}=\frac{36}9+\frac{25}9=\boxed{\frac{61}9}$\\[4pt]
$\frac34+\frac68=\frac68+\frac68=\frac{12}8=\boxed{\frac32}$\\[4pt]
$\frac56-\frac14=\frac{10}{12}-\frac3{12}=\boxed{\frac7{12}}$
\end{minipage}\hfill
\begin{minipage}[t]{0.49\linewidth}
Tres fraccions: \ $\frac45+\frac27-\frac12$\\ mcm 70: \ $\frac{56}{70}+\frac{20}{70}-\frac{35}{70}=\boxed{\frac{41}{70}}$\\[4pt]
$2-\frac35=\frac{10}5-\frac35=\boxed{\frac75}$
\end{minipage}
\end{exemple}

\vfill
\begin{exemple}[Exemple: \ dels deures, n'he fet $\frac13$ al matí i $\frac25$ a la vesprada. Quina part he fet? Quina me'n falta?]
Fet: \ $\frac13+\frac25=\frac5{15}+\frac6{15}=\boxed{\frac{11}{15}}$ \qquad Falta: \ $1-\frac{11}{15}=\frac{15}{15}-\frac{11}{15}=\boxed{\frac4{15}}$
\end{exemple}

\vfill
\begin{truc}[Simplifica al final]
Dona el resultat \textbf{irreductible}: \ $\frac{12}8=\frac32$.
\end{truc}

\vfill
% ============================================================ 4.9
\clearpage
\seccio{sec9}{Multiplicació i divisió de fraccions}{Multipliquem en línia; dividim en creu.}

\vfill
\begin{tcbraster}[raster columns=2,raster equal height,raster column skip=10pt]
\begin{tcolorbox}[fmbase,colback=fmK!8,colframe=fmK,coltitle=white,title={\faTimes\ \ Multiplicació},before upper=\centering]
{\Large$\dfrac ab\cdot\dfrac cd=\dfrac{a\cdot c}{b\cdot d}$}\par\smallskip
$\dfrac64\cdot\dfrac25=\dfrac{12}{20}=\dfrac35$
\end{tcolorbox}
\begin{tcolorbox}[fmbase,colback=fmA!8,colframe=fmA,coltitle=white,title={\faDivide\ \ Divisió},before upper=\centering]
{\Large$\dfrac ab:\dfrac cd=\dfrac{a\cdot d}{b\cdot c}$}\par\smallskip
$\dfrac43:\dfrac52=\dfrac{8}{15}$ \ \nota{en creu}
\end{tcolorbox}
\end{tcbraster}

\vfill
\begin{minipage}[t]{0.5\linewidth}
\begin{idea}[Fracció inversa]
La inversa de $\frac ab$ és $\frac ba$: \textbf{girem} la fracció.\\[3pt]
$\frac32\cdot\frac23=\frac66=1$ \qquad La inversa de 5 és $\frac15$.\\[3pt]
Dividir és \textbf{multiplicar per la inversa}. El 0 no té inversa.
\end{idea}
\end{minipage}\hfill
\begin{minipage}[t]{0.47\linewidth}
\begin{truc}[Simplifica abans]
Un numerador amb un denominador:\\[3pt]
$\dfrac{6}{4}\cdot\dfrac{2}{5}=\dfrac{3\cdot\cancel{2}\cdot\cancel{2}}{\cancel{2}\cdot\cancel{2}\cdot5}=\dfrac35$
\end{truc}
\end{minipage}

\vfill
\begin{tcbraster}[raster columns=2,raster equal height,raster column skip=10pt]
\begin{tcolorbox}[fmbase,colback=fmdefbg,colframe=fmdef,coltitle=white,title={Un natural per una fracció}]
Multipliquem només el numerador:\par\smallskip
{\centering$2\cdot\dfrac56=\dfrac{10}6=\dfrac53$\par}
\end{tcolorbox}
\begin{tcolorbox}[fmbase,colback=fmdefbg,colframe=fmdef,coltitle=white,title={La fracció d'una fracció}]
És un producte:\par\smallskip
{\centering$\dfrac23$ de $\dfrac35=\dfrac23\cdot\dfrac35=\dfrac6{15}=\dfrac25$\par}
\end{tcolorbox}
\end{tcbraster}

\vfill
\begin{exemple}[Exemples]
\begin{minipage}[t]{0.48\linewidth}
$\frac34\cdot\frac{8}{9}=\frac{24}{36}=\boxed{\frac23}$\\[4pt]
$5\cdot\frac{3}{10}=\frac{15}{10}=\boxed{\frac32}$\\[4pt]
$\frac25\cdot\frac52=\boxed{1}$ \ \nota{inverses}
\end{minipage}\hfill
\begin{minipage}[t]{0.49\linewidth}
$\frac56:\frac{10}{3}=\frac{15}{60}=\boxed{\frac14}$\\[4pt]
$\frac34:3=\frac34:\frac31=\boxed{\frac14}$\\[4pt]
$6:\frac23=\frac{18}{2}=\boxed{9}$ \ \nota{quants $\frac23$ caben en 6?}
\end{minipage}
\end{exemple}

\vfill
\begin{compte}[No cal comú denominador]
Per a multiplicar i dividir \textbf{no} reduïm a comú denominador: això només per a sumar i restar.
\end{compte}

\vfill
% ============================================================ 4.10
\clearpage
\seccio{sec10}{Operacions combinades}{La mateixa jerarquia que amb els enters.}

\vfill
\begin{minipage}[c]{0.4\linewidth}\raggedright
\pas{1} Parèntesis $\to$ \pas{2} potències $\to$ \pas{3} $\cdot$ i $:$ $\to$ \pas{4} $+$ i $-$.\\[8pt]
\faCheck\ \textbf{Simplifica} cada resultat parcial.\\[6pt]
\faCheck\ Comú denominador \textbf{només} per als termes que sumes en aquell pas.\\[6pt]
\faCheck\ Un $-$ davant d'un parèntesi afecta \textbf{tot} el resultat.
\end{minipage}\hfill
\begin{minipage}[c]{0.57\linewidth}
\begin{exemple}[Exemple: \ $\dfrac34-\dfrac12\cdot\left(\dfrac56-\dfrac13\right)$]
\vspace{-6pt}
\begin{alignat*}{2}
&=\dfrac34-\dfrac12\cdot\left(\dfrac56-\dfrac26\right)=\dfrac34-\dfrac12\cdot\dfrac12 &\quad& \text{\nota{\pas{1}}}\\
&=\dfrac34-\dfrac14 && \text{\nota{\pas{3}}}\\
&=\dfrac24=\boxed{\dfrac12} && \text{\nota{\pas{4}}}
\end{alignat*}
\end{exemple}
\end{minipage}

\vfill
\begin{exemple}[Exemple: \ $\left(1-\dfrac25\right)^{2}:\dfrac{9}{10}+\dfrac13$]
$=\left(\dfrac35\right)^{2}:\dfrac{9}{10}+\dfrac13=\dfrac{9}{25}:\dfrac{9}{10}+\dfrac13=\dfrac{90}{225}+\dfrac13=\dfrac25+\dfrac13=\dfrac{6}{15}+\dfrac{5}{15}=\boxed{\dfrac{11}{15}}$
\end{exemple}

\vfill
\begin{minipage}[t]{0.48\linewidth}
\begin{exemple}[Exemple: \ $\dfrac23+\dfrac12\cdot\dfrac45$]
$=\dfrac23+\dfrac{4}{10}$ \ \nota{\pas{3} primer}\\[3pt]
$=\dfrac23+\dfrac25=\dfrac{10}{15}+\dfrac{6}{15}=\boxed{\dfrac{16}{15}}$
\end{exemple}
\end{minipage}\hfill
\begin{minipage}[t]{0.49\linewidth}
\begin{exemple}[Exemple: \ $\left(\dfrac34-\dfrac12\right):\dfrac58$]
$=\dfrac14:\dfrac58$ \ \nota{\pas{1}}\\[3pt]
$=\dfrac{8}{20}=\boxed{\dfrac25}$ \ \nota{en creu}
\end{exemple}
\end{minipage}

\vfill
\begin{minipage}[t]{0.48\linewidth}
\begin{exemple}[Exemple: \ $\dfrac56-\dfrac23:4$]
$=\dfrac56-\dfrac2{12}=\dfrac56-\dfrac16$\\[3pt]
$=\dfrac46=\boxed{\dfrac23}$
\end{exemple}
\end{minipage}\hfill
\begin{minipage}[t]{0.49\linewidth}
\begin{exemple}[Exemple: \ $1+\dfrac35\cdot\left(2-\dfrac13\right)$]
$=1+\dfrac35\cdot\dfrac53$ \ \nota{\pas{1}}\\[3pt]
$=1+1=\boxed{2}$
\end{exemple}
\end{minipage}

\vfill
\begin{compte}[Primer el producte]
$\dfrac23+\dfrac12\cdot\dfrac45$ \ \textbf{no} és \ \mal{\left(\dfrac23+\dfrac12\right)\cdot\dfrac45}: sense parèntesis, la multiplicació va abans que la suma.
\end{compte}

\vfill
% ============================================================ 4.11
\clearpage
\seccio{sec11}{Problemes amb fraccions}{Un dibuix amb una barra ho aclareix tot.}

\vfill
\begin{tcbraster}[raster columns=2,raster equal height=rows,raster column skip=10pt]
\begin{tcolorbox}[fmbase,colback=fmdefbg,colframe=fmdef,coltitle=white,title={Fracció d'una quantitat}]
Dividim entre el denominador i multipliquem pel numerador.\par\smallskip
{\centering$\dfrac34$ de $200=\dfrac{3\cdot200}{4}=150$\par}
\end{tcolorbox}
\begin{tcolorbox}[fmbase,colback=fmdefbg,colframe=fmdef,coltitle=white,title={Fracció d'una fracció}]
Les \textbf{multipliquem}.\par\smallskip
{\centering$\dfrac23$ de $\dfrac35=\dfrac23\cdot\dfrac35=\dfrac25$\par}
\end{tcolorbox}
\begin{tcolorbox}[fmbase,colback=fmdefbg,colframe=fmdef,coltitle=white,title={La part que queda}]
El total és $1=\dfrac{b}{b}$. Si gastem $\dfrac{a}{b}$, en queda $1-\dfrac{a}{b}$.\par\smallskip
{\centering Gastem $\dfrac25$: \ en queda $\dfrac55-\dfrac25=\dfrac35$\par}
\end{tcolorbox}
\begin{tcolorbox}[fmbase,colback=fmdefbg,colframe=fmdef,coltitle=white,title={Problema invers}]
Si $\dfrac{a}{b}$ del total és $P$, el total és $P:\dfrac{a}{b}$.\par\smallskip
{\centering$\dfrac23$ del total són 40: \ $40:\dfrac23=60$\par}
\end{tcolorbox}
\end{tcbraster}

\vfill
\begin{exemple}[Exemple: \ l'Anna gasta $\frac25$ dels seus 60 € en roba i $\frac13$ del que li queda en un llibre. Quant li queda?]
\begin{minipage}[c]{0.45\linewidth}\centering
\begin{tikzpicture}[x=1.3cm]
  \foreach \i in {0,...,4} {\filldraw[fill=\ifnum\i<2 fmB!35\else white\fi,draw=fmB,line width=1pt] (\i,0) rectangle (\i+1,0.6);}
  \node[font=\small] at (1,0.3) {roba};
  \filldraw[fill=fmA!35,draw=fmA,line width=1pt] (2,0) rectangle (3,0.6); \node[font=\small] at (2.5,0.3) {llibre};
  \node[font=\small] at (4,0.3) {queda};
  \draw[decorate,decoration={brace,amplitude=5pt},line width=1pt] (0,0.7) -- (5,0.7) node[midway,above=6pt]{60 €};
  \foreach \i in {0,...,4} \node[font=\small,text=fmmuted] at (\i+0.5,-0.3) {12 €};
\end{tikzpicture}
\end{minipage}\hfill
\begin{minipage}[c]{0.52\linewidth}
Roba: \ $\dfrac25$ de $60=24$ € \ $\to$ \ en queden 36 €\\[4pt]
Llibre: \ $\dfrac13$ de $36=12$ €\\[4pt]
Li queden \ $36-12=\boxed{24\text{ €}}$
\end{minipage}
\end{exemple}

\vfill
\begin{exemple}[Exemple: \ un dipòsit està ple als $\frac23$ i té 40 litres. Quina capacitat té?]
\begin{minipage}[c]{0.4\linewidth}\centering
\begin{tikzpicture}[x=1.4cm]
  \foreach \i in {0,1,2} {\filldraw[fill=\ifnum\i<2 fmA!35\else white\fi,draw=fmA,line width=1pt] (\i,0) rectangle (\i+1,0.6);}
  \draw[decorate,decoration={brace,amplitude=5pt},line width=1pt] (0,0.7) -- (2,0.7) node[midway,above=6pt]{40 L};
  \foreach \i in {0,1,2} \node[font=\small,text=fmmuted] at (\i+0.5,-0.3) {20 L};
\end{tikzpicture}
\end{minipage}\hfill
\begin{minipage}[c]{0.57\linewidth}
Problema \textbf{invers}: \ $40:\dfrac23=\dfrac{40\cdot3}{2}=\boxed{60\text{ L}}$\\[4pt]
\nota{Cada terç són 20 L; els tres terços, 60 L.}
\end{minipage}
\end{exemple}

\vfill
\begin{minipage}[t]{0.5\linewidth}
\begin{compte}[De quina quantitat?]
El $\frac13$ és \textbf{del que queda} (36 €), no del total (60 €). Llig bé de què és cada fracció.
\end{compte}
\end{minipage}\hfill
\begin{minipage}[t]{0.47\linewidth}
\begin{truc}[Comprova que té sentit]
Respon amb una \textbf{frase} i amb les \textbf{unitats}: «Li queden 24 €.»
\end{truc}
\end{minipage}

\vfill
% ============================================================ resum
\resumtitol
\begin{tcbraster}[raster columns=2,raster equal height=rows,raster column skip=14pt,raster row skip=20pt,raster before skip=12pt]
\rtile{4.1}{Fracció}{$\frac ab$: $b$ parts iguals, en prenem $a$.\\[1ex] Pròpia $<1$; impròpia $>1$.}
\rtile{4.2}{D'un nombre}{$\frac35$ de 40 $=(40:5)\cdot3=24$.}
\rtile{4.3}{Percentatges}{$35\,\%=\frac{35}{100}$.\\[1ex] 20\,\% de 180 $=\frac{20\cdot180}{100}$.}
\rtile{4.4}{Equivalents}{Amplificar o simplificar.\\[1ex] Irreductible: entre el MCD.}
\rtile{4.5}{Comprovar}{$\frac ab=\frac cd\iff a\cdot d=b\cdot c$.}
\rtile{4.6}{Comú denominador}{El mcm dels denominadors.}
\rtile{4.7}{Comparar}{Mateix denominador:\\[1ex] el numerador més gran.}
\rtile{4.8}{Sumes i restes}{Comú denominador.\\[1ex] Sumem els numeradors.}
\rtile{4.9}{Producte i quocient}{Producte: en línia.\\[1ex] Quocient: en creu (per la inversa).}
\rtile{4.10}{Combinades}{Parèntesis, $\cdot$ i $:$, $+$ i $-$.\\[1ex] Simplifica.}
\rtile{4.11}{Problemes}{Fracció d'una quantitat\\[1ex] i de la part que queda.}
\end{tcbraster}
\end{document}
